\documentclass[11pt,a4paper]{article}
\usepackage[margin=2.3cm]{geometry}
\usepackage{amsmath,amssymb,bm}
\usepackage{booktabs}
\usepackage{enumitem}
\usepackage{xcolor}
\usepackage[hidelinks]{hyperref}
\usepackage{framed}

\newcommand{\R}{\mathbb{R}}
\newcommand{\E}{\mathbb{E}}
\newcommand{\Var}{\mathrm{Var}}
\newcommand{\N}{\mathcal{N}}
\newcommand{\btheta}{\bm{\theta}}
\newcommand{\bphi}{\bm{\phi}}
\newcommand{\by}{\mathbf{y}}
\newcommand{\br}{\mathbf{r}}
\newcommand{\bk}{\mathbf{k}}

\title{POPS-based Bayesian optimization with data-conditioned\\predictive entropy search}
\author{Algorithm notes for the \texttt{pops-bo} interactive demo\\\small\texttt{examples/pops\_bo\_web} in \url{https://github.com/tomswinburne/pops-bo}}
\date{\today}

\begin{document}
\maketitle

\begin{abstract}
These notes specify the algorithm implemented in \texttt{pops\_bo.js}. A structured but
misspecified parametric surrogate is equipped with a POPS ensemble (pointwise optimal parameter
sets) that represents model-form uncertainty. It is augmented with Bayesian-ridge epistemic draws, and
its complexity is chosen by Bayesian evidence. Experiments are selected by an approximate
predictive entropy search over the location of the optimum, following Lee \emph{et al.}~\cite{lee2026}.
We add two changes. First, the hypothetical outcomes are conditioned on data already measured near
the candidate through a discrepancy Gaussian process. Second, when the optimum location is resolved,
the method switches to the expected improvement of the best measured value. Under that criterion
replicate measurements have exactly zero value. Together these remove the repeated sampling of
the same point that the unconditioned procedure exhibits.
\end{abstract}

\section{Problem and notation}
We maximise an expensive, noisy response $y = f(x) + \epsilon$, with $\epsilon\sim\N(0,\sigma_n^2)$ and
$\sigma_n$ known, over a candidate grid $\{x_g\}_{g=1}^{G}\subset[0,1]$. After $N$ experiments the data are
$\mathcal D = \{(x_n,y_n)\}_{n=1}^N$. Performance is the \emph{best measured} regret
$f^\star - \max_n f(x_n)$ (noise-free $f$ at the measured locations); we also report the regret of the
recommendation $\hat x$.

\paragraph{Toy problem.} $f(x) = 0.3x + \sin(\omega x + 0.3) + \delta\sum_{k=2}^{8}\rho^{k-2}\sin(k\omega x + \phi_k)$,
with $\omega = 2\pi/P$, $\phi_k = 0.7 + 1.9(k-2)$, and defaults $P=0.45$, $\delta=0.3$, $\rho=0.6$, $\sigma_n=0.05$, $G=401$.
Setting $\rho=0$ gives the original single-harmonic specification. Every surrogate below is misspecified for $\delta>0$.

\section{Surrogate, prior and complexity}
The surrogate with complexity $K$ has $p = 2+2K$ features,
\begin{equation}
f_{\btheta}(x) = \btheta^\top\bphi_K(x),\qquad
\bphi_K(x) = \bigl(1,\,x,\,\sin\omega x,\,\cos\omega x,\,\dots,\,\sin K\omega x,\,\cos K\omega x\bigr)^\top .
\end{equation}
With prior $\btheta\sim\N(0,\tau^2 I)$ ($\tau=1$) and design matrix $\Phi\in\R^{N\times p}$, the posterior mode is the ridge solution
\begin{equation}
A = \Phi^\top\Phi + \lambda I,\quad \lambda = \sigma_n^2/\tau^2,\quad \btheta_* = A^{-1}\Phi^\top\by,\quad e_n = y_n - \bphi_n^\top\btheta_* .
\end{equation}

\paragraph{Model selection.} $K\in\{1,\dots,6\}$ is chosen at every step by maximising the log evidence
\begin{equation}
\log p(\by\mid K) = -\frac{\by^\top\by - \by^\top\Phi\btheta_*}{2\sigma_n^2} - \tfrac12\log|A| + \tfrac p2\log\lambda - \tfrac N2\log(2\pi\sigma_n^2),
\end{equation}
which equals $\log\N(\by;\,0,\,\sigma_n^2 I + \tau^2\Phi\Phi^\top)$. The Occam factor penalises unneeded harmonics. Misfit beyond the
noise penalises too few. The selected $K$ grows with $N$.

\section{Ensemble: POPS plus epistemic draws}
\paragraph{POPS members (Lee \emph{et al.}\ Eq.~6; Swinburne \& Perez~\cite{pops2025}).} For each observation $i$,
$\btheta_i = \arg\min_{\btheta} \|\Phi\btheta-\by\|^2 + \lambda\|\btheta\|^2$ subject to $\bphi_i^\top\btheta = y_i$. This has the
closed form (rank-one correction)
\begin{equation}
\btheta_i = \btheta_* + \frac{e_i}{h_i}\,A^{-1}\bphi_i,\qquad h_i = \bphi_i^\top A^{-1}\bphi_i .
\end{equation}
The approximation of Lee \emph{et al.}, a refit with $10\times$ loss weight on point $i$, is available and gives indistinguishable results.

\paragraph{Epistemic augmentation.} If $p\ge N$ the fit interpolates, $e_i\to 0$, and POPS alone would report a spuriously certain
complex model. Following \texttt{popsregression} (misspecification $+$ epistemic variance), the ensemble has
$M=\max(64,N)$ members
\begin{equation}
\btheta^{(m)} = \btheta_{i(m)} + \bm\xi_m,\qquad \bm\xi_m\sim\N(0,\sigma_n^2A^{-1}),\qquad i(m) = m \bmod N .
\end{equation}

\section{Conditioning outcomes on nearby data}
\label{sec:disc}
The approximate PES of Lee \emph{et al.}\ uses each member's prediction $f_{\btheta^{(k)}}(x)$ as a hypothetical outcome. Nothing in it
knows that $y$ has already been measured at, or near, $x$, so it has no anti-locality. We write the observation model with a
discrepancy term (Kennedy \& O'Hagan~\cite{ko2001}):
\begin{equation}
y(x) = f_{\btheta}(x) + d(x) + \epsilon,\qquad d\sim\mathcal{GP}\bigl(0,\,\tau_d^2\,e^{-(x-x')^2/2\ell^2}\bigr).
\end{equation}
Let $R_{nn'} = e^{-(x_n-x_{n'})^2/2\ell^2}$, $C = \tau_d^2R + \sigma_n^2 I$, and $\bk(x)_n = \tau_d^2e^{-(x-x_n)^2/2\ell^2}$.

\paragraph{Hyperparameters.} The discrepancy must absorb the misfit that POPS asserts. We therefore fit $(\ell,\tau_d^2)$ to the
residuals of the POPS members, $\br_i = \by - \Phi\btheta_i$, treated as independent draws, by grid search:
\begin{equation}
(\ell,\tau_d^2) = \arg\max \sum_{i=1}^{N}\log\N(\br_i;\,0,\,C).
\end{equation}
Fitting instead to the residuals of $\btheta_*$ alone gives $\tau_d\approx 0$ whenever the model fits within noise, and then no
conditioning takes place.

\paragraph{Conditioned member predictions.} For member $j$ with residuals $\br^{(j)} = \by - \Phi\btheta^{(j)}$:
\begin{align}
\mu_j(x) &= f_{\btheta^{(j)}}(x) + \bk(x)^\top C^{-1}\br^{(j)},\\
s^2_{\rm lat}(x) &= \tau_d^2 - \bk(x)^\top C^{-1}\bk(x),\qquad s^2(x) = s^2_{\rm lat}(x) + \sigma_n^2,\\
c(x,x') &= \tau_d^2 e^{-(x-x')^2/2\ell^2} - \bk(x)^\top C^{-1}\bk(x').
\end{align}
Near a measured $x_n$ every $\mu_j(x)$ is pulled to $y_n$ and $s^2$ shrinks to the noise. Far from data, $\mu_j\to f_{\btheta^{(j)}}$
and $s^2\to\tau_d^2+\sigma_n^2$, which recovers the paper's rule with likelihood width $\sigma^2_{\rm lik}\approx\tau_d^2+\sigma_n^2$.

\paragraph{Base weights.} With $i(j)$ the POPS parent of member $j$ and $c_i = |\{j: i(j)=i\}|$:
\begin{equation}
\log w^0_j = -\tfrac12\,\br_{i(j)}^\top C^{-1}\br_{i(j)} - \log c_{i(j)} + \text{const},\qquad \textstyle\sum_j w^0_j = 1 .
\end{equation}
Epistemic children share their parent's weight; re-weighting posterior draws by the data would double count.
Without conditioning the paper's form $-\|\br_{i(j)}\|^2/2\sigma^2_{\rm lik}$ is used, with $\sigma^2_{\rm lik} = \max(\mathrm{RSS}/(N-p),\sigma_n^2)$.

\section{Acquisition}
\paragraph{Predicted optima (Eq.~7).} $x^*_j = \arg\max_g \mu_j(x_g)$, and $H_0 = \tfrac12\log\bigl(2\pi e[\Var_{w^0}(x^*)+\varepsilon]\bigr)$,
with $\varepsilon = 0.01^2$, the resolution at which $x^*$ matters.

\paragraph{Data-conditioned approximate PES (Eqs.~8--12).} For candidate $x$ and each hypothetical outcome $y_k = \mu_k(x)$:
\begin{align}
\log w'_{j\mid k}(x) &= \log w^0_j - \frac{(\mu_j(x)-y_k)^2}{2s^2(x)}\quad\text{(normalised by log-sum-exp)},\\
H_k(x) &= \tfrac12\log\bigl(2\pi e\,[\Var_{w'_{\cdot\mid k}}(x^*)+\varepsilon]\bigr),\qquad
\alpha_{\rm PES}(x) = H_0 - \sum_k w^0_k H_k(x).
\end{align}
Only the weights change during acquisition; the ensemble is refitted after the real measurement. At a measured point all
$\mu_j$ coincide up to noise, so no member is reweighted and $\alpha_{\rm PES}\approx 0$.

\paragraph{Resolved optimum: improvement of the best measured value.} If $\max_x\alpha_{\rm PES}(x) < \gamma$ ($\gamma=0.05$ nats),
the members agree on $x^*$ to within $\varepsilon$ and PES has nothing left to resolve. The objective is the best \emph{measured}
value, so under hypothesis $j$ the gain from measuring $x$ is $\max\bigl(f_j(x)-\max_n f_j(x_n),\,0\bigr)$. With
$n_j = \arg\max_n\mu_j(x_n)$ and the Gaussian discrepancy,
\begin{equation}
\alpha_{\rm EI}(x) = \sum_j w^0_j\,\mathrm{EI}\bigl(\mu_j(x)-\mu_j(x_{n_j}),\;v_j(x)\bigr),\qquad
v_j(x) = s^2_{\rm lat}(x)+s^2_{\rm lat}(x_{n_j}) - 2c(x,x_{n_j}),
\end{equation}
where $\mathrm{EI}(d,v) = d\,\Phi(d/\sqrt v) + \sqrt v\,\varphi(d/\sqrt v)$. Because $f(x_m)\le\max_n f(x_n)$ for every measured $x_m$ and every
realisation, we set $\alpha_{\rm EI}(x_m)=0$ exactly. The variance of the \emph{difference} vanishes as $x\to x_{n_j}$, so
the criterion favours unmeasured points next to the peak and assigns no value to replicates. If $\alpha_{\rm EI}\equiv 0$ the
most informative candidate $\arg\max\alpha_{\rm PES}$ is measured.

\paragraph{Recommendation.} $\hat x = \arg\max_g\sum_j w^0_j\mu_j(x_g)$.

\begin{framed}
\noindent\textbf{Algorithm (one BO step).} Input $\mathcal D$, $\sigma_n$, grid.
\begin{enumerate}[leftmargin=1.6em,itemsep=1pt]
\item Choose $K$ by the evidence (Eq.~3); form $A$, $\btheta_*$, $e_n$, $h_n$.
\item POPS members $\btheta_i$ (Eq.~4); ensemble $\btheta^{(m)} = \btheta_{i(m)}+\bm\xi_m$ (Eq.~5).
\item Fit $(\ell,\tau_d^2)$ to the POPS residuals (Eq.~7); compute $\mu_j$, $s^2$, $s^2_{\rm lat}$, $c$ (Eqs.~8--10) and $w^0$ (Eq.~11).
\item $x^*_j$, $H_0$; evaluate $\alpha_{\rm PES}(x_g)$ for all $g$ (Eqs.~12--13).
\item If $\max\alpha_{\rm PES}\ge\gamma$: measure $\arg\max\alpha_{\rm PES}$; ties are broken by the highest ensemble mean.
\item Else: measure $\arg\max\alpha_{\rm EI}$ (Eq.~14, zero at measured points); if $\alpha_{\rm EI}\equiv0$, measure $\arg\max\alpha_{\rm PES}$.
\item Observe $y$, append to $\mathcal D$, refit from step 1.
\end{enumerate}
Cost per step is $O(GM^2)$ for PES, with $G=401$ and $M=64$, plus $O(N^3)$ grid fits; this is under 50\,ms in a browser.
\end{framed}

\section{Baselines and controls}
\begin{itemize}[itemsep=1pt]
\item \textbf{Greedy $\btheta_*$}: $\arg\max$ of the conditioned $\btheta_*$ prediction; the no-uncertainty ablation.
\item \textbf{POPS + UCB}: ensemble mean $+2\times$ ensemble sd; this uses response variance, not $x^*$ entropy.
\item \textbf{GP-RBF + EI}: stationary GP, with hyperparameters by marginal likelihood and noise $\ge\sigma_n^2/2$; noisy EI.
\item \textbf{GP-periodic + EI}: periodic kernel with known $P$ plus a linear term. It has the same structural prior without a fixed form.
\item \textbf{Paper-literal PES}: grid-scale $\varepsilon$, no conditioning, no exploitation rule.
\item \textbf{Minimum spacing}: an optional hard exclusion radius $r$. This was the earlier heuristic fix for repeats; $r=0$ by default.
\end{itemize}
All strategies share, per seed, the initial design (8 points, stratified in $[0,0.5]$) and the noise draws.

\section{Results}
\label{sec:results}
40 seeds, 15 iterations. These tables are generated from \texttt{results.json} by \texttt{docs/make\_tables.js}.
\noindent Settings: $\rho=0.6$, $\sigma_n=0.05$, 40 seeds, initial region $[0,0.5]$, $K$ by evidence, conditioning off, minimum spacing $r=0.005$.

\paragraph{Strategies, $P=0.45$, $\delta=0.3$.} Best measured regret (median, with mean) and recommendation regret.

\begin{center}\small\begin{tabular}{lrrrrr}\toprule
strategy & med.\ @3 & med.\ @5 & med.\ @15 & mean @15 & rec.\ @15\\\midrule
POPS + PES (conditioned) & 0.014 & 0.004 & 0.001 & 0.000 & 0.003\\
POPS + PES, paper-literal & 0.060 & 0.053 & 0.029 & 0.031 & 0.003\\
POPS + UCB & 0.008 & 0.001 & 0.001 & 0.001 & 0.001\\
Greedy $\theta_*$ & 0.017 & 0.001 & 0.000 & 0.000 & 0.003\\
GP-RBF + EI & 0.139 & 0.092 & 0.001 & 0.001 & 0.001\\
GP-periodic + EI & 0.003 & 0.001 & 0.000 & 0.000 & 0.001\\
\bottomrule\end{tabular}\end{center}

\paragraph{Strategies, $P=0.2$, $\delta=0.3$.} Best measured regret (median, with mean) and recommendation regret.

\begin{center}\small\begin{tabular}{lrrrrr}\toprule
strategy & med.\ @3 & med.\ @5 & med.\ @15 & mean @15 & rec.\ @15\\\midrule
POPS + PES (conditioned) & 0.013 & 0.003 & 0.000 & 0.001 & 0.003\\
POPS + PES, paper-literal & 0.179 & 0.132 & 0.057 & 0.053 & 0.007\\
POPS + UCB & 0.003 & 0.003 & 0.003 & 0.002 & 0.003\\
Greedy $\theta_*$ & 0.003 & 0.003 & 0.003 & 0.001 & 0.003\\
GP-RBF + EI & 0.142 & 0.130 & 0.003 & 0.006 & 0.005\\
GP-periodic + EI & 0.007 & 0.003 & 0.003 & 0.002 & 0.003\\
\bottomrule\end{tabular}\end{center}

\paragraph{Model complexity, $P=0.45$, $\delta=0.3$.} Best measured regret; the recommendation column shows the model bias.

\begin{center}\small\begin{tabular}{lrrrrr}\toprule
surrogate & med.\ @1 & med.\ @3 & med.\ @5 & med.\ @15 & rec.\ @15\\\midrule
POPS + PES, $K=1$ & 0.092 & 0.022 & 0.008 & 0.001 & 0.031\\
POPS + PES, $K=2$ & 0.076 & 0.008 & 0.001 & 0.001 & 0.008\\
POPS + PES, $K=4$ & 0.183 & 0.051 & 0.004 & 0.001 & 0.001\\
POPS + PES, $K=6$ & 0.273 & 0.125 & 0.048 & 0.001 & 0.003\\
POPS + PES, $K$ by evidence & 0.092 & 0.014 & 0.004 & 0.001 & 0.003\\
Greedy, $K$ by evidence & 0.092 & 0.017 & 0.001 & 0.000 & 0.003\\
GP-RBF + EI & 0.234 & 0.139 & 0.092 & 0.001 & 0.001\\
\bottomrule\end{tabular}\end{center}

\paragraph{Anti-locality ablation, $P=0.45$, $\delta=0.3$.} Fraction of BO steps that replicate a previous point (within $0.001$) or land within $0.005$ of one.

\begin{center}\small\begin{tabular}{lccrrrr}\toprule
strategy & conditioned & $r$ & replicates & within $0.005$ & med.\ @5 & mean @15\\\midrule
POPS + PES & no & 0 & 0\% & 72\% & 0.008 & 0.010\\
POPS + PES & no & 0.005 & 0\% & 0\% & 0.004 & 0.000\\
POPS + UCB & no & 0 & 0\% & 78\% & 0.004 & 0.013\\
POPS + UCB & no & 0.005 & 0\% & 0\% & 0.001 & 0.001\\
Greedy $\theta_*$ & no & 0 & 0\% & 86\% & 0.021 & 0.048\\
Greedy $\theta_*$ & no & 0.005 & 0\% & 0\% & 0.001 & 0.000\\
GP-RBF + EI & -- & 0 & 0\% & 38\% & 0.092 & 0.001\\
GP-RBF + EI & -- & 0.005 & 0\% & 0\% & 0.092 & 0.001\\
\bottomrule\end{tabular}\end{center}

\paragraph{Anti-locality ablation, $P=0.2$, $\delta=0.3$.} Fraction of BO steps that replicate a previous point (within $0.001$) or land within $0.005$ of one.

\begin{center}\small\begin{tabular}{lccrrrr}\toprule
strategy & conditioned & $r$ & replicates & within $0.005$ & med.\ @5 & mean @15\\\midrule
POPS + PES & no & 0 & 0\% & 73\% & 0.010 & 0.011\\
POPS + PES & no & 0.005 & 0\% & 0\% & 0.003 & 0.001\\
POPS + UCB & no & 0 & 0\% & 82\% & 0.007 & 0.007\\
POPS + UCB & no & 0.005 & 0\% & 0\% & 0.003 & 0.002\\
Greedy $\theta_*$ & no & 0 & 0\% & 89\% & 0.010 & 0.036\\
Greedy $\theta_*$ & no & 0.005 & 0\% & 0\% & 0.003 & 0.001\\
GP-RBF + EI & -- & 0 & 0\% & 7\% & 0.130 & 0.009\\
GP-RBF + EI & -- & 0.005 & 0\% & 0\% & 0.130 & 0.006\\
\bottomrule\end{tabular}\end{center}


\section{Limitations}
\begin{itemize}[itemsep=1pt]
\item The Gaussian entropy is crude for a multimodal $x^*$ distribution. Early on, disagreement between distant peaks dominates.
\item The conditioning adds a nonparametric component. With many local data the method approaches a GP with a parametric mean.
\item Discrepancy hyperparameters are point estimates on a grid.
\item If the initial data span less than one period, all members can agree on a wrong peak ($H_0\approx 0$).
\item The equation numbering for Lee \emph{et al.}\ follows the specification we were given; the preprint could not be retrieved while
these notes were written.
\end{itemize}

\begin{thebibliography}{9}
\bibitem{lee2026} J.~Lee, L.~J.~Augustine, P.~Yang, and D.~Perez,
\emph{Efficient Bayesian Optimization of Separation Chemistry in the Scarce-Data Regime Using Kinetic Reaction Networks},
ChemRxiv preprint (2026). \href{https://doi.org/10.26434/chemrxiv.15004837/v2}{doi:10.26434/chemrxiv.15004837/v2}.
\bibitem{pops2025} T.~D.~Swinburne and D.~Perez,
\emph{Parameter uncertainties for imperfect surrogate models in the low-noise regime},
Mach.\ Learn.: Sci.\ Technol.\ (2025). \href{https://doi.org/10.1088/2632-2153/ad9fce}{doi:10.1088/2632-2153/ad9fce}.
\bibitem{ko2001} M.~C.~Kennedy and A.~O'Hagan, \emph{Bayesian calibration of computer models}, J.\ R.\ Stat.\ Soc.\ B \textbf{63}, 425 (2001).
\bibitem{pes2014} J.~M.~Hern\'andez-Lobato, M.~W.~Hoffman, and Z.~Ghahramani, \emph{Predictive entropy search for efficient global optimization of black-box functions}, NeurIPS (2014).
\bibitem{jones1998} D.~R.~Jones, M.~Schonlau, and W.~J.~Welch, \emph{Efficient global optimization of expensive black-box functions}, J.\ Global Optim.\ \textbf{13}, 455 (1998).
\end{thebibliography}
\end{document}
